\chapter{Breastfeeding and Periods}
\index{Breastfeeding and Periods}

\section*{Definition and Overview}
This topic describes the relationship between breastfeeding and the menstrual cycle, particularly the delay in the return of periods after childbirth. The hormones that sustain milk production also suppress the reproductive cycle, so many breastfeeding women do not menstruate for months after giving birth. This pause, known as lactational amenorrhoea, varies considerably between women and from one pregnancy to the next. When periods do return, the first cycles are often irregular, heavy, or different from the pre-pregnancy pattern, and the timing is strongly influenced by how the baby is fed. Understanding how breastfeeding affects fertility also matters for family planning, because ovulation can resume before the first postpartum period appears, meaning breastfeeding alone is not a reliable method of contraception.

\section*{Epidemiology}
Almost all women experience some delay in the return of menstruation while breastfeeding. Women who breastfeed exclusively and frequently, including during the night, commonly remain amenorrhoeic for around six months, whereas women who formula-feed or mix-feed typically have their first period within two to three months of childbirth. Longer pauses are seen in women who breastfeed for extended periods and who feed on demand around the clock. When menstruation resumes, the first several cycles are frequently irregular, and ovulation may precede the first visible period in roughly half of breastfeeding women. This pattern is why relying on the absence of bleeding as a marker of infertility is unreliable and leads to unplanned pregnancies.

\section*{Aetiology and Pathophysiology}
The suppression of ovulation and menstruation during breastfeeding is a direct hormonal consequence of milk production.

\begin{itemize}
    \item \textbf{Prolactin:} the pituitary hormone that drives milk synthesis also inhibits the release of the gonadotrophins that normally stimulate the ovaries, keeping ovulation suppressed while feeding is frequent
    \item \textbf{Suckling stimulation:} the baby's sucking, especially at night, maintains elevated prolactin levels and sustained suppression of the reproductive axis
    \item \textbf{Feeding pattern:} exclusive, on-demand breastfeeding suppresses ovulation far more than infrequent or mixed feeding; introducing solids, long gaps between feeds, or night weaning allow the cycle to restart
    \item \textbf{Gradual hormonal recovery:} as feeding reduces, gonadotrophin levels recover and ovarian activity resumes, usually with ovulation occurring before the first visible period
    \item \textbf{Postpartum cycles:} early cycles are often anovulatory or irregular while the hypothalamic--pituitary--ovarian axis settles
\end{itemize}

\section*{Clinical Features}
\begin{itemize}
    \item \textbf{Lactational amenorrhoea:} the complete absence of periods for months while breastfeeding, which is normal
    \item \textbf{Return of periods:} the first period may be light, heavy, or prolonged, and may differ from the woman's usual pattern
    \item \textbf{Irregular cycles:} variable cycle length and timing over the first few cycles after return
    \item \textbf{Spotting:} light bleeding or brown discharge in the days before or instead of a period
    \item \textbf{Feeding-related changes:} a temporary dip in milk supply around the time a period is due, and sometimes increased breast tenderness that makes feeding less comfortable
    \item \textbf{Signs of returning ovulation:} mild pelvic pain (mittelschmerz), or changes in cervical mucus, before the first period
\end{itemize}

\section*{Differential Diagnosis}
\begin{itemize}
    \item \textbf{Lactational amenorrhoea:} the normal absence of periods while breastfeeding
    \item \textbf{New pregnancy:} missed, irregular, or unusual bleeding may reflect a new conception, particularly if contraception was not used
    \item \textbf{Postpartum bleeding:} early heavy or prolonged bleeding must be distinguished from the normal lochia of the weeks after birth
    \item \textbf{Thyroid or pituitary disorders:} conditions that disturb periods and can also affect milk supply
    \item \textbf{Polycystic ovary syndrome:} a common cause of irregular or absent periods in women of reproductive age
    \item \textbf{Uterine problems:} fibroids, polyps, or retained products of pregnancy causing abnormal bleeding
\end{itemize}

\section*{When to Seek Medical Care}
See a doctor if you are unsure whether bleeding is normal, if you think you may be pregnant, if periods have not returned within about a year of stopping breastfeeding, or if you have questions about contraception while breastfeeding. Heavy, painful, or unusual bleeding always warrants assessment.

\textbf{Contact your local emergency services for:}
\begin{itemize}
    \item Very heavy bleeding, or passing large clots, at any time after birth
    \item Bleeding accompanied by severe abdominal or pelvic pain
    \item Fever, chills, or a foul-smelling vaginal discharge suggesting infection
    \item Faintness, dizziness, or a racing heart in association with heavy bleeding
    \item A positive pregnancy test with pain or bleeding, which may indicate an ectopic pregnancy
\end{itemize}

\section*{Diagnosis and Investigations}
\begin{itemize}
    \item \textbf{History:} breastfeeding pattern, feeding frequency, night feeds, introduction of solids, contraception use, and associated symptoms
    \item \textbf{Pregnancy test:} to exclude a new pregnancy as the cause of missed or irregular periods
    \item \textbf{Examination:} abdominal and pelvic assessment if bleeding is heavy, painful, or associated with signs of infection
    \item \textbf{Blood tests:} thyroid function or reproductive hormone levels if periods do not return as expected after weaning, or if there are features of PCOS or other endocrine disorders
    \item \textbf{Pelvic ultrasound:} arranged when retained products of pregnancy, fibroids, polyps, or ovarian cysts are suspected
    \item \textbf{Full blood count:} to check for iron-deficiency anaemia when bleeding has been heavy or prolonged
\end{itemize}

\section*{Management}
\begin{itemize}
    \item \textbf{Reassurance:} in most cases, delayed or irregular periods while breastfeeding are normal and require no treatment
    \item \textbf{Contraception:} discussion of suitable methods; progestogen-only pills, injections, implants, intrauterine devices, and barrier methods are generally suitable during breastfeeding, while combined hormonal contraception requires individual review
    \item \textbf{Feeding patterns:} continued frequent, on-demand feeding, including at night, tends to delay the return of periods; introducing solids or reducing feeds brings them back sooner
    \item \textbf{Symptom relief:} simple analgesia and warmth for period-related breast tenderness, and comfortable feeding positions during sensitive days
    \item \textbf{Treating underlying causes:} if periods fail to return after weaning, or bleeding is heavy or irregular, treat the identified cause such as thyroid disease, PCOS, or uterine pathology
    \item \textbf{Follow-up:} medical review if periods remain absent after weaning or if bleeding is heavy, prolonged, or otherwise concerning
\end{itemize}

\section*{Complications}
\begin{itemize}
    \item \textbf{Unplanned pregnancy:} ovulation can resume before the first period, so breastfeeding is not a reliable contraceptive
    \item \textbf{Iron-deficiency anaemia:} from heavy or prolonged postpartum bleeding
    \item \textbf{Temporary milk supply dip:} a brief fall in supply around the time of a period, which usually recovers spontaneously
    \item \textbf{Anxiety about irregularity:} uncertainty about whether unusual cycles are normal
    \item \textbf{Late recognition of pregnancy:} if periods were assumed to be absent
    \item \textbf{Delayed diagnosis of pathology:} uncommon, but irregular bleeding may rarely be due to retained tissue, a polyp, or a fibroid
\end{itemize}

\section*{Prognosis and Outcomes}
For most women the return of periods during breastfeeding is a normal event that requires no treatment, and cycles settle into a regular pattern within a few months. The lactational pause tends to be longer with exclusive, frequent feeding and ends once feeding reduces or the baby begins solid foods. Because ovulation can return silently, reliable contraception is essential if another pregnancy is not desired. Once regular cycles resume, fertility returns to normal for the great majority of women, and any delay in periods after weaning usually resolves within a few months.

\section*{Prevention and Self-Care}
\begin{itemize}
    \item Use a reliable method of contraception from around three weeks after the birth if you do not wish to conceive again
    \item Do not rely on breastfeeding, or on the absence of periods, as a contraceptive method
    \item Track your cycle with an app or diary once periods return, to recognise patterns and flag concerns
    \item Report heavy, painful, or unusual bleeding to your doctor rather than dismissing it
    \item Continue to breastfeed on demand for as long as you wish; the choice is personal
    \item Seek medical advice if periods are slow to return after weaning, or are irregular or heavy
    \item Maintain adequate iron through diet, and consider iron testing if bleeding is heavy
\end{itemize}

\section*{Sources and Further Reading}
The following organisations provide reliable information on breastfeeding, the return of periods after birth, and contraception for breastfeeding women.

\begin{itemize}
    \item NHS. \textit{Breastfeeding, periods after birth, and contraception guidance.} https://www.nhs.uk/conditions/contraception/
    \item World Health Organization. \textit{Breastfeeding and reproductive health guidance.} https://www.who.int/
    \item Mayo Clinic. \textit{Return of periods after pregnancy and breastfeeding considerations.} https://www.mayoclinic.org/
    \item MedlinePlus. \textit{Breastfeeding and women's reproductive health information.} https://medlineplus.gov/breastfeeding.html
    \item Australian Breastfeeding Association. \textit{Breastfeeding support and information for families.} https://www.breastfeeding.asn.au/
    \item Family Planning Australia. \textit{Contraception options and reproductive health information.} https://www.fpnsw.org.au/
\end{itemize}
